\section{Additional Experiments}
\label{appendix_additional_experiments}

In this section, we present large-scale experiments of our CortexDebate and baseline methods to further demonstrate the superiority of CortexDebate. 

\textbf{Experimental Setup.} For each task (\textit{i.e.}, math, world knowledge question answering, reasoning, and long-context understanding), we conduct experiments on the more challenging one of the two datasets (\textit{i.e.}, MATH, MMLU-pro, GPQA, and LongBench). For each adopted dataset, we experiment on a subset of 1000 examples. Besides, the backbone models and evaluation metrics used in the experiments are the same as mentioned in Section~\ref{main_experimental_setup}.

\textbf{Results.} The experimental results are presented in Table~\ref{table_add_exp}. Consistent with the experimental results reported in Table~\ref{main_result_acc}, our proposed CortexDebate achieves the best performance on all the datasets compared with baseline methods. For instance, CortexDebate achieves a maximal RA of 56.30\% on MATH dataset, 58.90\% on MMLU-pro dataset, 36.60\% on GPQA dataset, and 59.63\% on LongBench dataset, respectively. Moreover, for each method, we calculate the average token numbers of the contexts input to one LLM agent on each dataset and present the results in Table~\ref{table_appendix_additional_exp_token}. Compared with the full debate methods (\textit{i.e.}, MLD, RECONCILE, ChatEval, and MPRC), our CortexDebate significantly reduces the length of the contextual input for each LLM agent, with a maximum reduction of 65.50\%. Moreover, compared with the part debate methods (\textit{i.e.}, GD and ND), our CortexDebate can reduce the input context length by at least 17.40\%.